\documentclass[a4paper]{article}     
\usepackage{amsfonts}
\usepackage{amsmath}
\usepackage{amssymb}
\usepackage{graphicx}
\usepackage{cancel}

\newcommand{\asa}{\Leftrightarrow} %als slechts als
\newcommand{\adR}{\Rightarrow} %als dan Rechts
\newcommand{\adL}{\Leftarrow}
\newcommand{\Lh}{\left(} %Links haakje
\newcommand{\Rh}{\right)}
\newcommand{\Lv}{\left[} %Links vierkanthaakje
\newcommand{\Rv}{\right]}
\newcommand{\La}{\left\{} %Links acolade
\newcommand{\Ra}{\right\}}
\newcommand{\Ras}{\right.} %Rechts acolade stelsel (omdat om stelsels te maken heb je een punt nodig
\newcommand{\pnR}{\rightarrow} %pijl naar Rechts
\newcommand{\limiet}{\lim\limits_}
\newcommand{\schrap}{\cancel}
\newcommand{\schrapb}{\bcancel} %backslash
\newcommand{\schrapk}{\xcancel} %kruis
\newcommand{\vervang}{\cancelto} 

\begin{document}

\noindent \large Auteur: Yoren Collier \\
\noindent \large Studierichting: Informatica\\
\noindent \large Oefeningenreeks 3B nr. 2.2\\

\medskip

\normalsize

$f(x) = \dfrac{1}{\sqrt{3 +2x}} = \Lh 3 +2x\Rh ^{\tfrac{-1}{2}}$\\

$\dfrac{df}{dx}(x) = \dfrac{-1}{\schrap{2}} \Lh 3 +2x \Rh ^{\tfrac{-3}{2}} \cdot \schrap{2}$\\

$\dfrac{d^2f}{dx^2}(x) = -1 \cdot \dfrac{-3}{\schrap{2}} \Lh 3 +2x \Rh ^{\tfrac{-5}{2}} \cdot \schrap{2}$\\

$\dfrac{d^3f}{dx^3}(x) = -1 \cdot -3 \cdot \dfrac{-5}{\schrap{2}} \Lh 3+2x \Rh ^{\tfrac{-7}{2}} \cdot \schrap{2} = \dfrac{-15}{\sqrt{\Lh 3 +2x \Rh ^7}}$\\

\begin{thebibliography}{999}
%\bibitem{a} Referentie
\end{thebibliography}
\end{document} 